\chapter{Conclusions}
\label{chap:conclusions}

This thesis started from a federated multi-agent reinforcement learning scheme that assigns channels and
transmit powers to \ac{V2V} links \cite{li+22fedmarl}. In that scheme, the vehicles decide on their own, but
they learn through the \ac{BS}: it averages their models and computes the reward from which all of them learn.
We implemented the scheme, proposed to take the aggregation away from the \ac{BS} and to give each agent its
own reward, and evaluated these proposals in 216 runs, with 5 to 30 \ac{V2V} pairs. This chapter answers the
research questions (Section~\ref{sec:conc-answers}), sums up the contributions and other findings
(Sections~\ref{sec:conc-contributions} and~\ref{sec:conc-findings}), discusses the limitations
(Section~\ref{sec:conc-limitations}), and proposes future work (Section~\ref{sec:conc-future}).

%%
\section{Answers to the research questions}
% Respuesta directa a PI1 (quitar la BS de la agregación) y a PI2 (recompensa individual), con los números principales.
\label{sec:conc-answers}

\paragraph{RQ1: Can the BS be removed from the aggregation of models without losing performance, and at what
cost?} Yes. With the individual reward, gossip among neighbors beats \ac{FedAvg} in all 12 paired comparisons
with 10 to 25 pairs, by 459 to 1754 points of reward per episode. Over all six values of $K$, it wins 14 of 18
comparisons ($p \approx 0.015$). With 5 pairs, the two tie. With 30 pairs, \ac{FedAvg} beats gossip with all
its neighbors in two of the three seeds, but gossip with at most two neighbors per round beats \ac{FedAvg} in
all three. With the shared reward, the two forms of aggregation tie: gossip wins 10 of 18 comparisons
($p \approx 0.41$).

The cost has two sides. In total, gossip with all its neighbors transfers more models than \ac{FedAvg} from 10
pairs on, up to 3.9 times more with 30 pairs, because each vehicle has more neighbors when there are more
vehicles. But the load is spread out: its busiest vehicle takes part in 36\% to 43\% of the transfers that the
\ac{BS} handles in \ac{FedAvg}, and these transfers do not go through the cellular network. With at most two
neighbors per round, the total never exceeds that of \ac{FedAvg}, and the load of the busiest vehicle stops
growing with the number of pairs. With 30 pairs, it is a tenth of the load of the \ac{BS}. This variant also
keeps the reward: it beats \ac{FedAvg} in 16 of 18 comparisons ($p \approx 7 \times 10^{-4}$).

\paragraph{RQ2: Do the agents learn better from an individual reward?} Yes, by a wide margin. We judged all the
agents by the shared reward. Even so, the agents that learn from their individual reward score higher than the
agents that learn from the shared reward itself in 50 of 54 paired comparisons
($p \approx 2 \times 10^{-11}$), including all 18 with gossip. With the methods that share their models, the
advantage tends to grow with the number of pairs, although not steadily. With the shared reward, the signal of
each agent gets lost among those of all the others, and with many pairs the agents learn to stay silent almost
all the time. With 25 pairs, the pairs of \ac{FedAvg} and gossip miss their reliability condition in 92\% to
93\% of the slots.

\paragraph{The two together.} The two proposals reinforce each other: gossip consistently beats \ac{FedAvg} only
when the agents learn from the individual reward. Together, they beat our implementation of the original scheme
in all 18 paired comparisons, for every number of pairs and every seed. With 15 pairs, the reward is 1.8 times
higher, and with 25 pairs, it goes from $-1682$ to 1545.

%%
\section{Contributions}
% Las tres contribuciones y lo que se logró con cada una.
\label{sec:conc-contributions}

\begin{enumerate}
  \item \textbf{A tested and reproducible implementation of the scheme of \citet{li+22fedmarl}}, with its
    environment, its agents, and three forms of aggregation. It has 191 automated tests, and the tables and
    figures of the results are generated directly from the logs of the runs. It reproduces part of the
    behavior of the published scheme (Section~\ref{sec:res-replication}).
  \item \textbf{A decentralized architecture for federated learning}, in which each vehicle averages its model
    with those of its neighbors within \SI{300}{\meter}, without the \ac{BS}. It also has a variant with a
    bounded cost, which uses at most $S$ neighbors per round.
  \item \textbf{An individual reward} that judges each agent by what its own action can change. Measured with
    the reward of all the agents, it gives better results than the shared reward.
\end{enumerate}

%%
\section{Other findings}
% Hallazgos secundarios: el factor de descuento importa, compartir con todos no siempre ayuda y la ventaja de gossip aparece después de la exploración.
\label{sec:conc-findings}

\paragraph{The discount factor matters.} \citet{li+22fedmarl} do not give the value of $\gamma$. With 0.99, the
value used by \citet{mnih+15dqn} and often taken as the default, our implementation loses 29\% of its best reward when
exploration ends and finishes below Greedy. With 0.9, it loses only 3\%. This result comes from a single seed
and 10 pairs.

\paragraph{Sharing with everyone does not always help.} With the individual reward, \ac{FedAvg} beats the
agents that learn alone in only 9 of 18 comparisons, and gossip beats them in 16 of 18. With 10 to 20 pairs,
averaging with all the agents gives the same result as not sharing, or a worse one. Whom an agent shares with
matters.

\paragraph{The advantage of gossip appears after exploration.} \ac{FedAvg} and gossip learn at the same pace
while the agents explore. When exploration ends, with the individual reward and 10 to 25 pairs, \ac{FedAvg}
drops and gossip keeps what it learned. With 30 pairs, this no longer holds. This thesis did not find out why
the two separate at that point.

%%
\section{Limitations}
% Simulación simplificada, implementación distinta del esquema publicado, recompensa individual que no es del todo local, intercambios ideales y demasiado frecuentes, evidencia estadística moderada y distancia a lo que es posible lograr.
\label{sec:conc-limitations}

\paragraph{A simplified simulation.} All the results come from a simulator with one cell, one path-loss law,
uniform positions that do not change within an episode, and constant traffic. With realistic mobility and
propagation, the size of the effects may change.

\paragraph{The implementation is not the published scheme.} It differs from \citet{li+22fedmarl} in several
training choices, and we tested only some of them with their values. It also does not reproduce their delay or,
with the shared reward, their fall in the rate of the \acp{CUE}. So the results show that the proposal beats this
implementation of the original scheme. They do not show that it would beat the published scheme with its own
settings.

\paragraph{The individual reward is not fully local.} Its term for the \ac{CUE} of the chosen channel needs the
\ac{BS} to broadcast the rate of its users in every slot. The reward also ignores the interference that a pair
causes to the other pairs on its channel. In this problem, that did not stop it from working better than the
shared reward, but it could matter more in denser scenarios. In addition, our results measure two changes
together: the terms that belong to each agent and the queue measured at the end of the slot.

\paragraph{Ideal and too frequent exchanges.} The models are transferred without losses, without delay, and
without using the spectrum that the agents share. Rounds every \SI{50}{\milli\second} would need about
\SI{19.5}{Mbps} per vehicle with \ac{FedAvg}. The comparison between the forms of aggregation holds in relative
terms, but it does not estimate the cost of a real system.

\paragraph{Moderate statistical evidence.} Three seeds per configuration give consistent paired comparisons, but
they cannot measure the size of the differences precisely. We singled out the range of 10 to 25 pairs and the
variant with $S = 2$ after seeing the results, so the comparisons over all six values of $K$ are the safer
numbers. We also tuned the configuration with a single seed, with \ac{FedAvg}, and with 10 pairs. See
Section~\ref{sec:res-threats} for the statistical caveats.

\paragraph{Far from what is possible.} No learning method reaches the threshold rule. The contributions improve
learning compared with our implementation of the original scheme, but they do not close that gap.

%%
\section{Future work}
% Escenarios más realistas, comparación con el esquema publicado, mejores recompensas, comunicación más barata y realista, entender por qué ayudan los vecinos, mejores entradas y entrenamiento, y robustez.
\label{sec:conc-future}

\paragraph{More realistic scenarios.} Evaluate the proposals with vehicles that move along roads, with the full
propagation models of the \ac{3GPP}, and with packet traffic, and measure how much the advantages of gossip and
of the individual reward change.

\paragraph{A closer comparison with the published scheme.} Test the learning rate and the learning frequency of
\citet{li+22fedmarl}, which the diagnosis did not try, to find out whether the published scheme has the same drop.

\paragraph{Better rewards.} Try difference rewards \cite{wolpert+01collectives,agogino+08rewards}, which also
count the harm that a pair causes to the other pairs on its channel and are still easy to learn from. Estimate
the effect on the \ac{CUE} locally, so that the agents do not depend on a broadcast from the \ac{BS}.

\paragraph{Cheaper and more realistic communication.} Reduce how often the rounds happen, compress the models,
and use the fact that a vehicle can send its model to all its neighbors with a single broadcast. Model the
losses and delays of the channel between vehicles. Study asynchronous gossip, in which each vehicle averages
when it can, without common rounds.

\paragraph{Why neighbors help.} Test the explanation of Section~\ref{sec:res-discussion}, that neighbors help
because they are in similar situations. One way is to measure the dispersion of the weights according to the
distance between agents. Another is to weigh each neighbor by how similar its situation is. Also find out why
the gap between \ac{FedAvg} and gossip appears only after exploration ends.

\paragraph{Better inputs and training.} Give the agents an estimate of the large-scale gain of their link as an
input, since today they must deduce it from 15 noisy gains. Look for ways to avoid the drop when exploration
ends.

\paragraph{Robustness.} Study what happens when vehicles enter and leave the cell, when exchanges fail, or when
a neighbor sends harmful models. Decentralized aggregation spreads this last risk, but it does not remove it.

\bigskip
\noindent In short, the \ac{BS} does not need to be at the center of learning. Vehicles can learn to share the
spectrum by averaging their models with their neighbors and by each learning from its own reward. In this
simulator, they learn better this way than with our implementation of the centralized scheme.
